% Part: many-valued-logic
% Chapter: three-valued-logics
% Section: lukasiewicz

\documentclass[../../../include/open-logic-section]{subfiles}

\begin{document}

\olfileid[id]{mvl}{inf}{luk}

\olsection{Logika \L ukasiewicz}

\begin{editorial}
  Ini adalah ``rintisan'' singkat untuk suatu bagian mengenai logika \L
  ukasiewicz bernilai tak hingga.
\end{editorial}

\begin{defn}\ollabel{def:lukasiewicz} Logika \L ukasiewicz bernilai tak
hingga~$\LogLuk[\infty]$ didefinisikan dengan menggunakan matriks berikut:
\begin{enumerate}
  \item Bahasa proposisional baku~$\Lang L_0$ dengan
  $\lnot$, $\land$, $\lor$, $\lif$.
  \item Himpunan nilai kebenaran~$V_\infty$.
  \item $1$ merupakan satu-satunya nilai kebenaran yang ditetapkan, yakni
  $V^+ = \{1\}$.
  \item Fungsi kebenaran diberikan oleh fungsi-fungsi berikut:
  \begin{align*}
    \tf{\lnot}[\LogLuk](x) & = 1 - x\\
    \tf{\land}[\LogLuk](x,y) & = \min(x,y)\\
    \tf{\lor}[\LogLuk](x,y) & = \max(x,y)\\
    \tf{\lif}[\LogLuk](x,y) & = \min(1,1-(x-y)) = \begin{cases}
      1 & \text{jika } x \le y\\
      1-(x-y) & \text{selain itu.}
    \end{cases}
    \end{align*}
\end{enumerate}
Logika \L ukasiewicz bernilai~$m$ didefinisikan dengan cara yang sama,
kecuali $V = V_m$.
\end{defn}

\begin{prop}
  Logika $\LogLuk[3]$ yang didefinisikan oleh
  \olref[thr][luk]{def:lukasiewicz} sama dengan $\LogLuk[3]$ yang
  didefinisikan oleh \olref{def:lukasiewicz}.
\end{prop}

\begin{proof}
  Hal ini dapat dilihat dengan membandingkan tabel kebenaran untuk
  konektif-konektif yang diberikan dalam
  \olref[thr][luk]{def:lukasiewicz} dengan tabel kebenaran yang ditentukan
  oleh persamaan-persamaan dalam \olref{def:lukasiewicz}:
  \begin{center}
    \begin{tabular}{c|c} 
      $\tf{\lnot}$ & \\ 
      \hline  
      $1$ & $0$ \\ 
      $1/2$ & $1/2$ \\
      $0$ & $1$ 
    \end{tabular}
    \quad
    \begin{tabular}{c|ccc} 
      $\tf{\land}[\LogLuk[3]]$ & $1$ & $1/2$ & $0$ \\ 
      \hline 
      $1$ & $1$ & $1/2$ & $0$ \\ 
      $1/2$ & $1/2$ & $1/2$ & $0$\\ 
      $0$ & $0$ & $0$ & $0$ 
    \end{tabular}
    \\[2ex]
    \begin{tabular}{c|ccc} 
      $\tf{\lor}[\LogLuk[3]]$ & $1$ & $1/2$ & $0$ \\ 
      \hline 
      $1$ & $1$ & $1$ & $1$ \\ 
      $1/2$ & $1$ & $1/2$ & $1/2$ \\
      $0$ & $1$ & $1/2$ & $0$ 
    \end{tabular}
    \quad
    \begin{tabular}{c|ccc} 
      $\tf{\lif}[\LogLuk[3]]$ & $1$ & $1/2$ & $0$ \\ 
      \hline 
      $1$ & $1$ & $1/2$ & $0$ \\ 
      $1/2$ & $1$ & $1$ & $1/2$  \\ 
      $0$ & $1$ & $1$ & $1$ 
    \end{tabular}
  \end{center} 
\end{proof}

\begin{prop}\ollabel{prop:luk-infty-m}
  Jika $\Gamma \Entails[\LogLuk[\infty]] !B$, maka $\Gamma
  \Entails[\LogLuk[m]] !B$ untuk setiap~$m \ge 2$.
\end{prop}

\begin{proof}
  Latihan.
\end{proof}

\begin{prob}
  Buktikan \olref[mvl][inf][luk]{prop:luk-infty-m}.
\end{prob}

Faktanya, kebalikannya juga berlaku.

Logika \L ukasiewicz bernilai tak hingga merupakan logika kabur yang paling
populer. Dalam literatur logika kabur, implikasi sering didefinisikan sebagai
$\lnot !A \lor !B$. Hasilnya adalah logika Kleene kuat bernilai tak hingga.

\begin{prob}
  Tunjukkan bahwa $(p \lif q) \lor (q \lif p)$ merupakan tautologi
  dari~$\LogLuk[\infty]$.
\end{prob}

\end{document}
